
Neural networks trained on datasets with spurious correlations exhibit systematic failures on minority subgroups. On the Waterbirds dataset, ResNet-50 achieves 97.2\% average accuracy but only 72.6\% worst-group accuracy \citep{Sagawa2020Distributionally}. Spurious correlations --- statistical associations between input features and labels that hold in training data but break under distribution shift --- cause models to learn shortcuts rather than core features \citep{Geirhos2020Shortcut}.

Existing work has established that deep neural networks exhibit simplicity bias, preferentially learning simpler decision boundaries that exploit spurious correlations \citep{Geirhos2020Shortcut}. Group Distributionally Robust Optimization (Group DRO) addresses this by minimizing worst-group loss rather than average loss, improving Waterbirds worst-group accuracy from 72.6\% to 91.4\% \citep{Sagawa2020Distributionally}. Temporal learning dynamics studies show that networks learn simple, often spurious, features earlier than complex core features \citep{Toneva2019Empirical}. However, systematic comparison of how architectural components --- specifically normalization layers --- affect spurious correlation learning dynamics has not been conducted.

This work investigates whether Batch Normalization and Layer Normalization exhibit different worst-group accuracy gaps when trained on spurious correlation tasks. We hypothesize that Batch Normalization's batch-level statistics amplify spurious correlations present in training batches, while Layer Normalization's instance-level normalization does not. To test this, we compare ResNet-18 with Batch Normalization (ResNet-18-BN) and ResNet-18 with Layer Normalization (ResNet-18-LN) at matched average accuracy checkpoints to eliminate training speed confounds.

We make three contributions:

\begin{enumerate}
\item \textbf{Quantitative gap difference}: ResNet-18-BN exhibits 9.41 percentage points higher worst-group accuracy gap than ResNet-18-LN when both reach 90\% average accuracy on synthetic spurious correlation data (p < 0.001, Cohen's d = 3.94, 10/10 seeds).
\end{enumerate}

\begin{enumerate}
\item \textbf{Gradient mechanism}: Batch Normalization shows 26.23\% higher gradient flow toward spurious-aligned samples during epochs 0--19 compared to Layer Normalization (p < 0.001, Cohen's d = 4.32).
\end{enumerate}

\begin{enumerate}
\item \textbf{Accuracy-matched comparison methodology}: Comparing architectures at matched average accuracy checkpoints (e.g., 90\%) eliminates training speed confounds while revealing temporal dynamics.
\end{enumerate}

These findings are based on proof-of-concept experiments with synthetic data due to WILDS server unavailability. Real dataset validation is required before claiming generalization to real-world spurious correlation tasks.
